\documentclass[11pt]{article}
% Shared print layout. The web's provisional design and MathJax are independent.
\usepackage[utf8]{inputenc}
\usepackage[T1]{fontenc}
\usepackage[spanish,es-nodecimaldot]{babel}
\usepackage{lmodern}
\usepackage{amsmath,amssymb}
\usepackage{graphicx}
\usepackage{geometry}
\geometry{a4paper,left=20mm,right=20mm,top=25mm,bottom=22mm,headheight=16pt,headsep=8mm}
\usepackage{fancyhdr,lastpage}
\usepackage{enumitem}
\usepackage{xcolor}
\usepackage{pgfplots}
\pgfplotsset{compat=1.18}
\usepgfplotslibrary{fillbetween}
\usepackage[hidelinks,pdfencoding=auto]{hyperref}
\hypersetup{pdflang=es-ES,pdfauthor={},pdfsubject={Matemáticas Aplicadas a las Ciencias Sociales II}}
\pagestyle{fancy}
\fancyhf{}
\fancyhead[L]{\small Matemáticas Aplicadas a las Ciencias Sociales II}
\fancyhead[R]{\small 2.\textsuperscript{o} Bachillerato}
\fancyfoot[L]{\footnotesize Distribuciones II | Edición revisada 1.0}
\fancyfoot[R]{\footnotesize Página \thepage\ de \pageref{LastPage}}
\renewcommand{\headrulewidth}{0.4pt}
\setlength{\parindent}{0pt}
\setlength{\parskip}{0.5em}
\setlist{itemsep=0.3em,topsep=0.3em}
\newcommand{\apartado}[1]{\subsection*{#1}}
\pgfplotsset{
  revision/.style={
    width=0.94\linewidth,height=7.5cm,
    axis lines=middle,grid=major,grid style={gray!18},
    tick label style={font=\small},label style={font=\small},
    xlabel={$x$},ylabel={$y$},samples=160,
    legend style={at={(0.5,-0.16)},anchor=north,draw=none,font=\small},
    clip=true
  }
}

\hypersetup{pdftitle={Distribuciones II - Soluciones revisadas}}
\begin{document}
\section*{Ejercicio 1: bicicletas y binomial}
Soluciones revisadas a partir de las incluidas en F002. Los datos y la puntuación del examen se conservan.

Se emplea el modelo del original: ensayos independientes y probabilidad constante de uso $p=1/5$. Si existiera muestreo sin reemplazo de una población pequeña conocida, habría que revisar esa hipótesis; el enunciado no da ese tamaño.
\[X\sim B(8,0{,}2),\qquad P(X=k)=\binom8k(0{,}2)^k(0{,}8)^{8-k}.\]
\subsection*{(a) Ninguno usa bicicleta}
\[P(X=0)=(0{,}8)^8=\boxed{0{,}16777216}.\]
\subsection*{(b) Al menos seis}
\begin{align*}
P(X=6)&=28(0{,}2)^6(0{,}8)^2=0{,}00114688,\\
P(X=7)&=8(0{,}2)^7(0{,}8)=0{,}00008192,\\
P(X=8)&=(0{,}2)^8=0{,}00000256.
\end{align*}
\[\boxed{P(X\geq6)=0{,}00123136}\quad(0{,}123136\%).\]
\textbf{Corrección M02:} el original escribe $(0{,}2)^8=0{,}0000256$, diez veces el valor correcto. Además suma términos redondeados antes del final. Se conservan aquí todos los decimales exactos de esta binomial.
\subsection*{(c) Al menos uno}
\[P(X\geq1)=1-P(X=0)=\boxed{0{,}83222784}.\]
Todas las probabilidades están entre 0 y 1 y la última es el complemento de la primera.
\clearpage
\section*{Ejercicio 2: normal y bonificación}
Se conserva la convención del original $N(\mu,\sigma)$, con $\mu=70$ y $\sigma=15$. Sea $Z=(X-70)/15\sim N(0,1)$ y $\Phi$ su función de distribución acumulada.
\subsection*{(a) Entre 60 y 90}
\[
P(60\leq X\leq90)=\Phi(4/3)-\Phi(-2/3)\approx0{,}65629624.
\]
El \textbf{porcentaje} es aproximadamente $\boxed{65{,}63\%}$.

La tabla del original redondea primero los valores tipificados a 1,33 y -0,67 y obtiene 0,6568, es decir, 65,68\%. Es una aproximación de tabla; la diferencia se debe al redondeo de los valores tipificados. No se presenta como igualdad exacta.
\subsection*{(b) Por debajo de 50}
\[P(X<50)=\Phi(-4/3)\approx\boxed{0{,}09121122}.\]
Con tabla y $z\approx-1{,}33$ se obtiene aproximadamente 0,0918, como en el original. El cálculo sin redondeo intermedio da 0,0912 a cuatro decimales.
\subsection*{(c) El 18\% superior}
\[P(X\geq x)=0{,}18\quad\Longleftrightarrow\quad\Phi\left(\frac{x-70}{15}\right)=0{,}82.\]
Por tanto,
\[z=\Phi^{-1}(0{,}82)\approx0{,}91536509,\qquad
\boxed{x\approx83{,}73048\text{ puntos}}.\]
A dos decimales, aproximadamente 83,73 puntos. Si una aplicación exigiera puntuaciones enteras, tendría que especificarlo; no se añade esa condición al examen.
\clearpage
\section*{Ejercicio 3: productos exitosos}
Bajo el modelo de resultados independientes con probabilidad constante,
\[X\sim B(100,0{,}6),\quad\mu=60,\quad\sigma=\sqrt{24}.\]
El suceso pedido es $X\geq70$, incluidos exactamente 70 éxitos.
\subsection*{Cálculo binomial exacto}
\[P(X\geq70)=\sum_{k=70}^{100}\binom{100}{k}(0{,}6)^k(0{,}4)^{100-k}
\approx\boxed{0{,}02478282}.\]
Es aproximadamente el 2,4783\%.
\subsection*{Aproximación normal conservada del original}
Los valores $np=60$ y $n(1-p)=40$ permiten una aproximación normal razonable. Se usa una variable auxiliar $Y\sim N(60,\sqrt{24})$; $X$ sigue siendo binomial.

Con \textbf{corrección de continuidad},
\begin{align*}
P(X\geq70)&\approx P(Y\geq69{,}5),\\
z&=\frac{69{,}5-60}{\sqrt{24}}\approx1{,}93917938,\\
P(Y\geq69{,}5)&\approx\boxed{0{,}02623975}.
\end{align*}
Con tabla y $z\approx1{,}94$, resulta 0,0262, el valor aproximado del original. No se confunde con la probabilidad binomial exacta. La diferencia absoluta entre los cálculos sin redondeo es aproximadamente 0,00145693.

Se redondea al final y se identifica el método utilizado. No basta afirmar que cualquier muestra con $n$ grande admite automáticamente una buena aproximación: importan también $p$ y las hipótesis del modelo.
\end{document}
